% !TeX program = xelatex
% ============================================================================
%  队友章节：视频 matting 与应用   ——   负责人：张竣翔
% ----------------------------------------------------------------------------
%  · 本文件由主文件 slides.tex 通过 % !TeX program = xelatex
% ============================================================================
%  队友章节：视频 matting 与应用   ——   负责人：张竣翔
% ----------------------------------------------------------------------------
%  · 本文件由主文件 slides.tex 通过 % !TeX program = xelatex
% ============================================================================
%  队友章节：视频 matting 与应用   ——   负责人：张竣翔
% ----------------------------------------------------------------------------
%  · 本文件由主文件 slides.tex 通过 % !TeX program = xelatex
% ============================================================================
%  队友章节：视频 matting 与应用   ——   负责人：张竣翔
% ----------------------------------------------------------------------------
%  · 本文件由主文件 slides.tex 通过 \input{sections/video.tex} 引入，
%    继承主文件导言区的全部宏 / 配色 / 字体，无需自带 \documentclass 或 \begin{document}。
%  · 直接在下面增删 \begin{frame}[t]{标题} ... \end{frame} 即可。
%  · 可复用样式：
%       \hilite{} \tcard{} \ccard{} \notestrip{} \metricchip{} \mpill[] \pill[]
%     可选 accent：nankai、accentblue、softgreen、softorange、softgray、deepred
% ============================================================================

\section{视频与应用}

\begin{frame}[t]{视频抠图与时序平滑}
  \scriptsize
  \linespread{0.98}\selectfont
  
  \begin{columns}[T,onlytextwidth]
    \column{0.48\textwidth}
      \begin{block}{问题：视频 alpha 抖动}
        \begin{itemize}
          \setlength{\itemsep}{0.2em}
          \item MatAnyone2~{\scriptsize[7]} 逐帧推理时，alpha 存在帧间\hilite{抖动与边界闪烁}；
          \item 直接替换背景会放大这些不稳定，导致视觉不自然；
          \item 需要一种\hilite{保持细节}的时序平滑策略。
        \end{itemize}
      \end{block}
      
      \vspace{0.2em}
      \begin{block}{解决方案：边界带时序平滑}
        \centering
        \resizebox{0.95\linewidth}{!}{%
        \begin{tikzpicture}[
          >={Latex[length=2.5mm,width=2mm]},
          node distance=3mm,
          every node/.style={font=\scriptsize},
          vcard/.style={rounded corners=2pt, line width=0.8pt, align=center, inner sep=3pt},
          vblue/.style={vcard, draw=accentblue!82!black, fill=accentblue!12},
          vpurple/.style={vcard, draw=nankai!88!black, fill=nankai!12},
          vgreen/.style={vcard, draw=softgreen!72!black, fill=softgreen!14},
          varrow/.style={-{Latex[length=2.5mm,width=2mm]}, line width=1.2pt, draw=nankai!80!black}
        ]
          \node[vblue] (input) {\textbf{输入}\\alpha video};
          \node[vpurple, right=of input] (detect) {\textbf{边界检测}\\$0.05<\alpha<0.95$};
          \node[vgreen, right=of detect] (smooth) {\textbf{时序平滑}\\EMA / Bidirectional};
          \node[vblue, right=of smooth] (output) {\textbf{输出}\\稳定 alpha};
          \draw[varrow] (input) -- (detect);
          \draw[varrow] (detect) -- (smooth);
          \draw[varrow] (smooth) -- (output);
        \end{tikzpicture}}
        
        \vspace{0.2em}
        \scriptsize
        仅在边界区域（$\alpha\in[0.05,0.95]$）施加平滑，\hilite{保护前景核区与背景}，避免过度模糊。
      \end{block}
      
    \column{0.48\textwidth}
      \begin{block}{7 种平滑策略对比}
        \centering
        \includegraphics[width=0.95\linewidth]{pic/figure6_performance_ranking.png}
        \vspace{0.1em}
        {\scriptsize TC 提升排名（左）与 SAD 变化排名（右）}
      \end{block}
  \end{columns}
\end{frame}

\begin{frame}[t]{7 种平滑策略性能评估}
  \scriptsize
  \linespread{0.50}\selectfont
  
  \begin{block}{实验设置}
    \scriptsize
    \begin{itemize}
      \setlength{\itemsep}{0.1em}
      \item \textbf{数据集：}YouTubeMatte (10 视频) + VideoMatt (10 视频)，共 20 个测试视频；
      \item \textbf{评估指标：}TC（时序一致性）、SAD（绝对误差和）、MSE、Gradient；
      \item \textbf{方法：}EMA、Gaussian、Median、Adaptive、Guided EMA、Bidirectional、Flow。
    \end{itemize}
  \end{block}
  
  \vspace{0.1em}
  \centering
  \resizebox{0.9\linewidth}{!}{%
  \begin{tabular}{lcccc}
    \toprule
    \textbf{方法} & \textbf{TC 提升 (\%)} & \textbf{SAD 变化 (\%)} & \textbf{MSE 变化 (\%)} & \textbf{细节保留} \\
    \midrule
    Median & 6.68 & +8.88 & +23.91 & Poor \\
    Gaussian & 5.50 & +5.40 & +11.42 & Fair \\
    Adaptive & 3.62 & +3.12 & +5.88 & Fair \\
    Guided EMA & 3.49 & +2.99 & +5.28 & Fair \\
    \textbf{Bidirectional} & \textbf{2.89} & \textbf{+0.44} & \textbf{-0.12} & \textbf{Excellent} \\
    EMA & 2.59 & +1.64 & +2.40 & Good \\
    Flow & 1.55 & +2.15 & +3.62 & Fair \\
    \bottomrule
  \end{tabular}}
  
  
  \vspace{0.1em}
  \notestrip{nankai}{\scriptsize
  \textbf{关键发现：}Bidirectional 是\hilite{唯一能同时改善 TC 和保持精度}的方法，细节保留最佳。
  }
\end{frame}

\begin{frame}[t]{方法对比散点图}
  \scriptsize
  \linespread{0.92}\selectfont
  
  \begin{columns}[T,onlytextwidth]
    \column{0.60\textwidth}
      \centering
      \includegraphics[width=\linewidth]{pic/figure1_tc_vs_sad_scatter.png}
      \vspace{0.1em}
      {\scriptsize  每个点代表一个视频的最佳结果}
      
    \column{0.4\textwidth}
      \begin{block}{解读}
        \begin{itemize}
          \setlength{\itemsep}{0.2em}
          \item 绿色区域：理想方向（TC↑，SAD↓）
          \item \textbf{Bidirectional}：最接近理想区域
          \item Median：TC 提升最大，但 SAD 增加显著
        \end{itemize}
      \end{block}
      
      \vspace{0.3em}
      \begin{block}{最佳参数}
        \scriptsize
        \begin{tabular}{ll}
          \toprule
          \textbf{方法} & \textbf{最佳参数} \\
          \midrule
          Median & window=7 \\
          Gaussian & sigma=2.0, window=7 \\
          \textbf{Bidirectional} & \textbf{alpha=0.7} \\
          \bottomrule
        \end{tabular}
      \end{block}
  \end{columns}
\end{frame}

\begin{frame}[t]{各方法性能分布箱线图}
  \scriptsize
  \linespread{0.98}\selectfont
  
  \begin{columns}[T,onlytextwidth]
    \column{0.7\textwidth}
      \centering
      \includegraphics[width=\linewidth]{pic/figure2_method_comparison_boxplot.png}
      \vspace{0.1em}
      {\scriptsize 左：TC 提升分布 \quad | \quad 右：SAD 变化分布}
      
    \column{0.32\textwidth}
      \begin{block}{箱线图解读}
        \begin{itemize}
          \setlength{\itemsep}{0.2em}
          \item 箱体 = 四分位距（50\% 数据）
          \item 横线 = 中位数
          \item 圆形 = 异常值
        \end{itemize}
      \end{block}
      
      \vspace{0.3em}
      \begin{block}{关键观察}
        \begin{itemize}
          \setlength{\itemsep}{0.2em}
          \item \textbf{Bidirectional}：SAD 分布最集中，异常值最少
          \item \textbf{Median}：TC 提升高但波动大
          \item \textbf{Gaussian}：TC 提升分布较稳定
        \end{itemize}
      \end{block}
  \end{columns}
  
  \vspace{0.1em}
  \notestrip{nankai}{\scriptsize
  \textbf{结论：}Bidirectional 的 SAD 变化分布最集中，说明其精度表现\hilite{最稳定}，不同视频间差异最小。
  }
\end{frame}

\begin{frame}[t]{Bidirectional 平滑效果对比}
  \scriptsize
  \linespread{0.98}\selectfont
  
  \begin{columns}[T,onlytextwidth]
    \column{0.55\textwidth}
      \centering
      \includegraphics[width=0.98\linewidth]{pic/gif_frames/frame_0001.png}
      \vspace{0.1em}
      {\tiny 左：原始 alpha \quad | \quad 右：Bidirectional 平滑后}
      
      \vspace{0.2em}
      \href{run:pic/youtubematte_motion_0016_bidirectional_comparison_slow_hq.gif}{%
        \color{accentblue}{\underline{[Click] 点击播放动画效果 (GIF)}}%
      }
      
      \vspace{0.3em}
      \begin{block}{视觉效果}
        \begin{itemize}
          \setlength{\itemsep}{0.2em}
          \item 发丝细节\hilite{无明显模糊}
          \item 边缘抖动\hilite{明显减弱}
          \item 背景替换后\hilite{闪烁减少}
        \end{itemize}
      \end{block}
      
    \column{0.42\textwidth}
      \begin{block}{性能指标}
        \scriptsize
        \begin{tabular}{lcc}
          \toprule
          \textbf{指标} & \textbf{数值} & \textbf{评价} \\
          \midrule
          TC 提升 & 2.89\% & 提升 \\
          SAD 变化 & +0.44\% & 几乎无损 \\
          MSE 变化 & -0.12\% & 改善 \\
          \bottomrule
        \end{tabular}
      \end{block}
      
      \vspace{0.3em}
      \begin{block}{优劣分析}
        \scriptsize
        \begin{itemize}
          \setlength{\itemsep}{0.1em}
          \item \textbf{优点：}精度几乎无损，发丝细节保留完整
          \item \textbf{缺点：}TC 提升幅度相对较小
        \end{itemize}
      \end{block}
  \end{columns}
  
  \vspace{0.2em}
  \notestrip{softgreen}{\scriptsize
  \textbf{结论：}Bidirectional 在\hilite{保持视觉质量的同时}实现时序一致性提升，是精度无损场景的最佳选择。
  }
\end{frame}

\begin{frame}[t]{Median 平滑效果对比}
  \scriptsize
  \linespread{0.98}\selectfont
  
  \begin{columns}[T,onlytextwidth]
    \column{0.55\textwidth}
      \centering
      \includegraphics[width=0.98\linewidth]{pic/median_frames/median_frame_0001.png}
      \vspace{0.1em}
      {\tiny 左：原始 alpha \quad | \quad 右：Median 平滑后}
      
      \vspace{0.2em}
      \href{run:pic/youtubematte_motion_0016_comparison_slow_hq.gif}{%
        \color{accentblue}{\underline{[Click] 点击播放动画效果 (GIF)}}%
      }
      
      \vspace{0.3em}
      \begin{block}{视觉效果}
        \begin{itemize}
          \setlength{\itemsep}{0.2em}
          \item 发丝细节\hilite{明显模糊}
          \item 边缘抖动\hilite{减少但细节丢失}
          \item 背景替换后\hilite{毛边明显}
        \end{itemize}
      \end{block}
      
    \column{0.42\textwidth}
      \begin{block}{性能指标}
        \scriptsize
        \begin{tabular}{lcc}
          \toprule
          \textbf{指标} & \textbf{数值} & \textbf{评价} \\
          \midrule
          TC 提升 & 6.68\% & 提升最大 \\
          SAD 变化 & +8.88\% & 恶化严重 \\
          MSE 变化 & +23.91\% & 恶化严重 \\
          \bottomrule
        \end{tabular}
      \end{block}
      
      \vspace{0.3em}
      \begin{block}{优劣分析}
        \scriptsize
        \begin{itemize}
          \setlength{\itemsep}{0.1em}
          \item \textbf{优点：}TC 提升最大（6.68\%），时序稳定性最好
          \item \textbf{缺点：}发丝模糊严重，视觉损失不可接受
        \end{itemize}
      \end{block}
  \end{columns}
  
  \vspace{0.2em}
  \notestrip{softorange}{\scriptsize
  \textbf{结论：}Median 数值最优但\hilite{视觉损失严重}，不适合需要精细边缘的实际应用。
  }
\end{frame}


\begin{frame}[t]{Gaussian 平滑效果对比}
  \scriptsize
  \linespread{0.98}\selectfont
  
  \begin{columns}[T,onlytextwidth]
    \column{0.55\textwidth}
      \centering
      
\includegraphics[width=0.98\linewidth]{pic/gaussian_frames/gaussian_frame_0001.png}
      \vspace{0.1em}
      {\tiny 左：原始 alpha \quad | \quad 右：Gaussian 平滑后 (sigma=2.0)}
      
      \vspace{0.2em}
      \href{run:pic/youtubematte_motion_0016_gaussian_comparison_slow_hq.gif}{%
        \color{accentblue}{\underline{[Click] 点击播放动画效果 (GIF)}}%
      }
      
      \vspace{0.3em}
      \begin{block}{视觉效果}
        \begin{itemize}
          \setlength{\itemsep}{0.2em}
          \item 发丝细节\hilite{轻微模糊}
          \item 边缘抖动\hilite{明显减少}
          \item 整体\hilite{平衡性较好}
        \end{itemize}
      \end{block}
      
    \column{0.42\textwidth}
      \begin{block}{性能指标}
        \scriptsize
        \begin{tabular}{lcc}
          \toprule
          \textbf{指标} & \textbf{数值} & \textbf{评价} \\
          \midrule
          TC 提升 & 4.0\% & 提升显著 \\
          SAD 变化 & +2.4\% & 轻微恶化 \\
          MSE 变化 & +3.0\% & 轻微恶化 \\
          \bottomrule
        \end{tabular}
      \end{block}
      
      \vspace{0.3em}
      \begin{block}{优劣分析}
        \scriptsize
        \begin{itemize}
          \setlength{\itemsep}{0.1em}
          \item \textbf{优点：}TC 提升显著（4.0\%），是 TC 和精度的最佳平衡点
          \item \textbf{缺点：}发丝有轻微模糊，精度略有下降
        \end{itemize}
      \end{block}
  \end{columns}
  
  \vspace{0.2em}
  \notestrip{accentblue}{\scriptsize
  \textbf{建议：}追求平衡可选 Gaussian (sigma=2.0, window=5)；追求精度无损可选 Bidirectional (alpha=0.7)。
  }
\end{frame}


\begin{frame}[t]{两个数据集上的性能对比}
  \scriptsize
  \linespread{0.98}\selectfont
  
  \begin{columns}[T,onlytextwidth]
    \column{0.7\textwidth}
      \centering
      \includegraphics[width=0.98\linewidth]{pic/figure3_dataset_comparison.png}
      \vspace{0.1em}
      {\scriptsize 左：TC 提升对比 \quad | \quad 右：SAD 变化对比}
      
    \column{0.32\textwidth}
      \begin{block}{关键观察}
        \begin{itemize}
          \setlength{\itemsep}{0.2em}
          \item YouTubeMatte 上效果\hilite{普遍更好}
          \item VideoMatt 更具挑战性
          \item \textbf{Bidirectional} 表现\hilite{最稳定}
        \end{itemize}
      \end{block}
      
      \vspace{0.5em}
      \begin{block}{差异分析}
        \scriptsize
        VideoMatt 包含更多\hilite{复杂运动}和\hilite{精细结构}，平滑难度更高。
      \end{block}
    \end{columns}
   \vspace{0.1em}
  \notestrip{accentblue}{\scriptsize
  \textbf{结论：}Bidirectional 在两个数据集上表现最稳定，适合不同场景的视频抠图任务。
  }
\end{frame}

\begin{frame}[t]{Gaussian 参数敏感性分析}
  \vspace{-0.5cm}
  \scriptsize
  \linespread{0.98}\selectfont
  
  \begin{columns}[T,onlytextwidth]
    \column{0.65\textwidth}
      \centering
      \includegraphics[width=0.98\linewidth]{pic/figure4_gaussian_parameter_sensitivity.png}
      \vspace{0.1em}
      {\scriptsize 气泡大小 = TC 提升，颜色 = SAD 变化}
      
    \column{0.32\textwidth}
      \begin{block}{参数趋势}
        \begin{itemize}
          \setlength{\itemsep}{0.2em}
          \item 增大 sigma → TC 提升增加
          \item 增大 sigma → SAD 恶化加剧
          \item 增大窗口 → 效果饱和
        \end{itemize}
      \end{block}
      
      \vspace{0.3em}
      \begin{block}{推荐参数}
        \centering
        \metricchip{softgreen}{sigma=2.0}{window=5} \\
        \vspace{0.1em}
        {\scriptsize TC 提升: 4.0\%, SAD 增加: 2.4\%}
      \end{block}
  \end{columns}
  
  \vspace{0.1em}
  \notestrip{accentblue}{\scriptsize
  \textbf{建议：}追求平衡可选 Gaussian；追求精度无损可选 Bidirectional。
  }
\end{frame}

\begin{frame}[t]{综合性能雷达图}
  \vspace{-0.5cm}
  \scriptsize
  \linespread{0.98}\selectfont
  
  \begin{columns}[T,onlytextwidth]
    \column{0.60\textwidth}
      \centering
      \includegraphics[width=0.95\linewidth]{pic/figure5_radar_chart.png}
      
    \column{0.32\textwidth}
      \begin{block}{解读}
        \begin{itemize}
          \setlength{\itemsep}{0.2em}
          \item 雷达图中 SAD/MSE/Gradient 已反向归一化，\hilite{值越大 = 性能越好}。
          \item \textbf{Bidirectional} 最均衡
          \item Median 偏科严重
        \end{itemize}
      \end{block}
      
      \vspace{0.3em}
      \begin{block}{总结}
        \centering
        \textbf{Bidirectional (alpha=0.7)} \\
        \vspace{0.1em}
        → 雷达图面积最大 \\
        → 推荐为默认平滑策略
      \end{block}
  \end{columns}
  
  
\end{frame}

\begin{frame}[t]{视频换背景 Demo}
  \scriptsize
  \linespread{0.98}\selectfont
  
  \begin{columns}[T,onlytextwidth]
    \column{0.48\textwidth}
      \begin{block}{平滑后 alpha 对比}
        \centering
        \includegraphics[width=0.95\linewidth]{pic/gif_frames/frame_0001.png}
        \vspace{0.1em}
        {\scriptsize 左：原始 \quad | \quad 右：Bidirectional 平滑后}
        
        \vspace{0.2em}
        \href{run:pic/youtubematte_motion_0016_bidirectional_slow_hq.gif}{%
          \color{accentblue}{\underline{[Click] 播放动画效果 (GIF)}}%
        }
      \end{block}
      
    \column{0.48\textwidth}
      \begin{block}{换背景效果}
        \centering
        
\includegraphics[width=0.35\linewidth]{pic/background.jpg}
        \vspace{0.1em}
        \\
        {\scriptsize 平滑后 alpha → RGBA → 背景替换}
        
        \vspace{0.2em}
        \href{run:pic/youtubematte_motion_0016_bidirectional_composer.mp4}{%
          \color{accentblue}{\underline{[Click] 播放换背景动画 (mp4)}}%
        }
      \end{block}
  \end{columns}
  
  \vspace{0.2em}
  \begin{block}{接口说明}
    \scriptsize
    \textbf{输入：}图像主线首帧 alpha / mask
    \textbf{处理：}MatAnyone2 + Bidirectional 时序平滑 (alpha=0.7)
    \textbf{输出：}逐帧稳定 alpha + 换背景合成视频
  \end{block}
  
  \vspace{0.1em}
  \notestrip{softgreen}{\scriptsize
  \textbf{效果：}Bidirectional 平滑后，背景替换视频的\hilite{边界闪烁明显减少}，同时\hilite{发丝细节保持清晰}。
  }
\end{frame} 引入，
%    继承主文件导言区的全部宏 / 配色 / 字体，无需自带 \documentclass 或 \begin{document}。
%  · 直接在下面增删 \begin{frame}[t]{标题} ... \end{frame} 即可。
%  · 可复用样式：
%       \hilite{} \tcard{} \ccard{} \notestrip{} \metricchip{} \mpill[] \pill[]
%     可选 accent：nankai、accentblue、softgreen、softorange、softgray、deepred
% ============================================================================

\section{视频与应用}

\begin{frame}[t]{视频抠图与时序平滑}
  \scriptsize
  \linespread{0.98}\selectfont
  
  \begin{columns}[T,onlytextwidth]
    \column{0.48\textwidth}
      \begin{block}{问题：视频 alpha 抖动}
        \begin{itemize}
          \setlength{\itemsep}{0.2em}
          \item MatAnyone2~{\scriptsize[7]} 逐帧推理时，alpha 存在帧间\hilite{抖动与边界闪烁}；
          \item 直接替换背景会放大这些不稳定，导致视觉不自然；
          \item 需要一种\hilite{保持细节}的时序平滑策略。
        \end{itemize}
      \end{block}
      
      \vspace{0.2em}
      \begin{block}{解决方案：边界带时序平滑}
        \centering
        \resizebox{0.95\linewidth}{!}{%
        \begin{tikzpicture}[
          >={Latex[length=2.5mm,width=2mm]},
          node distance=3mm,
          every node/.style={font=\scriptsize},
          vcard/.style={rounded corners=2pt, line width=0.8pt, align=center, inner sep=3pt},
          vblue/.style={vcard, draw=accentblue!82!black, fill=accentblue!12},
          vpurple/.style={vcard, draw=nankai!88!black, fill=nankai!12},
          vgreen/.style={vcard, draw=softgreen!72!black, fill=softgreen!14},
          varrow/.style={-{Latex[length=2.5mm,width=2mm]}, line width=1.2pt, draw=nankai!80!black}
        ]
          \node[vblue] (input) {\textbf{输入}\\alpha video};
          \node[vpurple, right=of input] (detect) {\textbf{边界检测}\\$0.05<\alpha<0.95$};
          \node[vgreen, right=of detect] (smooth) {\textbf{时序平滑}\\EMA / Bidirectional};
          \node[vblue, right=of smooth] (output) {\textbf{输出}\\稳定 alpha};
          \draw[varrow] (input) -- (detect);
          \draw[varrow] (detect) -- (smooth);
          \draw[varrow] (smooth) -- (output);
        \end{tikzpicture}}
        
        \vspace{0.2em}
        \scriptsize
        仅在边界区域（$\alpha\in[0.05,0.95]$）施加平滑，\hilite{保护前景核区与背景}，避免过度模糊。
      \end{block}
      
    \column{0.48\textwidth}
      \begin{block}{7 种平滑策略对比}
        \centering
        \includegraphics[width=0.95\linewidth]{pic/figure6_performance_ranking.png}
        \vspace{0.1em}
        {\scriptsize TC 提升排名（左）与 SAD 变化排名（右）}
      \end{block}
  \end{columns}
\end{frame}

\begin{frame}[t]{7 种平滑策略性能评估}
  \scriptsize
  \linespread{0.50}\selectfont
  
  \begin{block}{实验设置}
    \scriptsize
    \begin{itemize}
      \setlength{\itemsep}{0.1em}
      \item \textbf{数据集：}YouTubeMatte (10 视频) + VideoMatt (10 视频)，共 20 个测试视频；
      \item \textbf{评估指标：}TC（时序一致性）、SAD（绝对误差和）、MSE、Gradient；
      \item \textbf{方法：}EMA、Gaussian、Median、Adaptive、Guided EMA、Bidirectional、Flow。
    \end{itemize}
  \end{block}
  
  \vspace{0.1em}
  \centering
  \resizebox{0.9\linewidth}{!}{%
  \begin{tabular}{lcccc}
    \toprule
    \textbf{方法} & \textbf{TC 提升 (\%)} & \textbf{SAD 变化 (\%)} & \textbf{MSE 变化 (\%)} & \textbf{细节保留} \\
    \midrule
    Median & 6.68 & +8.88 & +23.91 & Poor \\
    Gaussian & 5.50 & +5.40 & +11.42 & Fair \\
    Adaptive & 3.62 & +3.12 & +5.88 & Fair \\
    Guided EMA & 3.49 & +2.99 & +5.28 & Fair \\
    \textbf{Bidirectional} & \textbf{2.89} & \textbf{+0.44} & \textbf{-0.12} & \textbf{Excellent} \\
    EMA & 2.59 & +1.64 & +2.40 & Good \\
    Flow & 1.55 & +2.15 & +3.62 & Fair \\
    \bottomrule
  \end{tabular}}
  
  
  \vspace{0.1em}
  \notestrip{nankai}{\scriptsize
  \textbf{关键发现：}Bidirectional 是\hilite{唯一能同时改善 TC 和保持精度}的方法，细节保留最佳。
  }
\end{frame}

\begin{frame}[t]{方法对比散点图}
  \scriptsize
  \linespread{0.92}\selectfont
  
  \begin{columns}[T,onlytextwidth]
    \column{0.60\textwidth}
      \centering
      \includegraphics[width=\linewidth]{pic/figure1_tc_vs_sad_scatter.png}
      \vspace{0.1em}
      {\scriptsize  每个点代表一个视频的最佳结果}
      
    \column{0.4\textwidth}
      \begin{block}{解读}
        \begin{itemize}
          \setlength{\itemsep}{0.2em}
          \item 绿色区域：理想方向（TC↑，SAD↓）
          \item \textbf{Bidirectional}：最接近理想区域
          \item Median：TC 提升最大，但 SAD 增加显著
        \end{itemize}
      \end{block}
      
      \vspace{0.3em}
      \begin{block}{最佳参数}
        \scriptsize
        \begin{tabular}{ll}
          \toprule
          \textbf{方法} & \textbf{最佳参数} \\
          \midrule
          Median & window=7 \\
          Gaussian & sigma=2.0, window=7 \\
          \textbf{Bidirectional} & \textbf{alpha=0.7} \\
          \bottomrule
        \end{tabular}
      \end{block}
  \end{columns}
\end{frame}

\begin{frame}[t]{各方法性能分布箱线图}
  \scriptsize
  \linespread{0.98}\selectfont
  
  \begin{columns}[T,onlytextwidth]
    \column{0.7\textwidth}
      \centering
      \includegraphics[width=\linewidth]{pic/figure2_method_comparison_boxplot.png}
      \vspace{0.1em}
      {\scriptsize 左：TC 提升分布 \quad | \quad 右：SAD 变化分布}
      
    \column{0.32\textwidth}
      \begin{block}{箱线图解读}
        \begin{itemize}
          \setlength{\itemsep}{0.2em}
          \item 箱体 = 四分位距（50\% 数据）
          \item 横线 = 中位数
          \item 圆形 = 异常值
        \end{itemize}
      \end{block}
      
      \vspace{0.3em}
      \begin{block}{关键观察}
        \begin{itemize}
          \setlength{\itemsep}{0.2em}
          \item \textbf{Bidirectional}：SAD 分布最集中，异常值最少
          \item \textbf{Median}：TC 提升高但波动大
          \item \textbf{Gaussian}：TC 提升分布较稳定
        \end{itemize}
      \end{block}
  \end{columns}
  
  \vspace{0.1em}
  \notestrip{nankai}{\scriptsize
  \textbf{结论：}Bidirectional 的 SAD 变化分布最集中，说明其精度表现\hilite{最稳定}，不同视频间差异最小。
  }
\end{frame}

\begin{frame}[t]{Bidirectional 平滑效果对比}
  \scriptsize
  \linespread{0.98}\selectfont
  
  \begin{columns}[T,onlytextwidth]
    \column{0.55\textwidth}
      \centering
      \includegraphics[width=0.98\linewidth]{pic/gif_frames/frame_0001.png}
      \vspace{0.1em}
      {\tiny 左：原始 alpha \quad | \quad 右：Bidirectional 平滑后}
      
      \vspace{0.2em}
      \href{run:pic/youtubematte_motion_0016_bidirectional_comparison_slow_hq.gif}{%
        \color{accentblue}{\underline{[Click] 点击播放动画效果 (GIF)}}%
      }
      
      \vspace{0.3em}
      \begin{block}{视觉效果}
        \begin{itemize}
          \setlength{\itemsep}{0.2em}
          \item 发丝细节\hilite{无明显模糊}
          \item 边缘抖动\hilite{明显减弱}
          \item 背景替换后\hilite{闪烁减少}
        \end{itemize}
      \end{block}
      
    \column{0.42\textwidth}
      \begin{block}{性能指标}
        \scriptsize
        \begin{tabular}{lcc}
          \toprule
          \textbf{指标} & \textbf{数值} & \textbf{评价} \\
          \midrule
          TC 提升 & 2.89\% & 提升 \\
          SAD 变化 & +0.44\% & 几乎无损 \\
          MSE 变化 & -0.12\% & 改善 \\
          \bottomrule
        \end{tabular}
      \end{block}
      
      \vspace{0.3em}
      \begin{block}{优劣分析}
        \scriptsize
        \begin{itemize}
          \setlength{\itemsep}{0.1em}
          \item \textbf{优点：}精度几乎无损，发丝细节保留完整
          \item \textbf{缺点：}TC 提升幅度相对较小
        \end{itemize}
      \end{block}
  \end{columns}
  
  \vspace{0.2em}
  \notestrip{softgreen}{\scriptsize
  \textbf{结论：}Bidirectional 在\hilite{保持视觉质量的同时}实现时序一致性提升，是精度无损场景的最佳选择。
  }
\end{frame}

\begin{frame}[t]{Median 平滑效果对比}
  \scriptsize
  \linespread{0.98}\selectfont
  
  \begin{columns}[T,onlytextwidth]
    \column{0.55\textwidth}
      \centering
      \includegraphics[width=0.98\linewidth]{pic/median_frames/median_frame_0001.png}
      \vspace{0.1em}
      {\tiny 左：原始 alpha \quad | \quad 右：Median 平滑后}
      
      \vspace{0.2em}
      \href{run:pic/youtubematte_motion_0016_comparison_slow_hq.gif}{%
        \color{accentblue}{\underline{[Click] 点击播放动画效果 (GIF)}}%
      }
      
      \vspace{0.3em}
      \begin{block}{视觉效果}
        \begin{itemize}
          \setlength{\itemsep}{0.2em}
          \item 发丝细节\hilite{明显模糊}
          \item 边缘抖动\hilite{减少但细节丢失}
          \item 背景替换后\hilite{毛边明显}
        \end{itemize}
      \end{block}
      
    \column{0.42\textwidth}
      \begin{block}{性能指标}
        \scriptsize
        \begin{tabular}{lcc}
          \toprule
          \textbf{指标} & \textbf{数值} & \textbf{评价} \\
          \midrule
          TC 提升 & 6.68\% & 提升最大 \\
          SAD 变化 & +8.88\% & 恶化严重 \\
          MSE 变化 & +23.91\% & 恶化严重 \\
          \bottomrule
        \end{tabular}
      \end{block}
      
      \vspace{0.3em}
      \begin{block}{优劣分析}
        \scriptsize
        \begin{itemize}
          \setlength{\itemsep}{0.1em}
          \item \textbf{优点：}TC 提升最大（6.68\%），时序稳定性最好
          \item \textbf{缺点：}发丝模糊严重，视觉损失不可接受
        \end{itemize}
      \end{block}
  \end{columns}
  
  \vspace{0.2em}
  \notestrip{softorange}{\scriptsize
  \textbf{结论：}Median 数值最优但\hilite{视觉损失严重}，不适合需要精细边缘的实际应用。
  }
\end{frame}


\begin{frame}[t]{Gaussian 平滑效果对比}
  \scriptsize
  \linespread{0.98}\selectfont
  
  \begin{columns}[T,onlytextwidth]
    \column{0.55\textwidth}
      \centering
      
\includegraphics[width=0.98\linewidth]{pic/gaussian_frames/gaussian_frame_0001.png}
      \vspace{0.1em}
      {\tiny 左：原始 alpha \quad | \quad 右：Gaussian 平滑后 (sigma=2.0)}
      
      \vspace{0.2em}
      \href{run:pic/youtubematte_motion_0016_gaussian_comparison_slow_hq.gif}{%
        \color{accentblue}{\underline{[Click] 点击播放动画效果 (GIF)}}%
      }
      
      \vspace{0.3em}
      \begin{block}{视觉效果}
        \begin{itemize}
          \setlength{\itemsep}{0.2em}
          \item 发丝细节\hilite{轻微模糊}
          \item 边缘抖动\hilite{明显减少}
          \item 整体\hilite{平衡性较好}
        \end{itemize}
      \end{block}
      
    \column{0.42\textwidth}
      \begin{block}{性能指标}
        \scriptsize
        \begin{tabular}{lcc}
          \toprule
          \textbf{指标} & \textbf{数值} & \textbf{评价} \\
          \midrule
          TC 提升 & 4.0\% & 提升显著 \\
          SAD 变化 & +2.4\% & 轻微恶化 \\
          MSE 变化 & +3.0\% & 轻微恶化 \\
          \bottomrule
        \end{tabular}
      \end{block}
      
      \vspace{0.3em}
      \begin{block}{优劣分析}
        \scriptsize
        \begin{itemize}
          \setlength{\itemsep}{0.1em}
          \item \textbf{优点：}TC 提升显著（4.0\%），是 TC 和精度的最佳平衡点
          \item \textbf{缺点：}发丝有轻微模糊，精度略有下降
        \end{itemize}
      \end{block}
  \end{columns}
  
  \vspace{0.2em}
  \notestrip{accentblue}{\scriptsize
  \textbf{建议：}追求平衡可选 Gaussian (sigma=2.0, window=5)；追求精度无损可选 Bidirectional (alpha=0.7)。
  }
\end{frame}


\begin{frame}[t]{两个数据集上的性能对比}
  \scriptsize
  \linespread{0.98}\selectfont
  
  \begin{columns}[T,onlytextwidth]
    \column{0.7\textwidth}
      \centering
      \includegraphics[width=0.98\linewidth]{pic/figure3_dataset_comparison.png}
      \vspace{0.1em}
      {\scriptsize 左：TC 提升对比 \quad | \quad 右：SAD 变化对比}
      
    \column{0.32\textwidth}
      \begin{block}{关键观察}
        \begin{itemize}
          \setlength{\itemsep}{0.2em}
          \item YouTubeMatte 上效果\hilite{普遍更好}
          \item VideoMatt 更具挑战性
          \item \textbf{Bidirectional} 表现\hilite{最稳定}
        \end{itemize}
      \end{block}
      
      \vspace{0.5em}
      \begin{block}{差异分析}
        \scriptsize
        VideoMatt 包含更多\hilite{复杂运动}和\hilite{精细结构}，平滑难度更高。
      \end{block}
    \end{columns}
   \vspace{0.1em}
  \notestrip{accentblue}{\scriptsize
  \textbf{结论：}Bidirectional 在两个数据集上表现最稳定，适合不同场景的视频抠图任务。
  }
\end{frame}

\begin{frame}[t]{Gaussian 参数敏感性分析}
  \vspace{-0.5cm}
  \scriptsize
  \linespread{0.98}\selectfont
  
  \begin{columns}[T,onlytextwidth]
    \column{0.65\textwidth}
      \centering
      \includegraphics[width=0.98\linewidth]{pic/figure4_gaussian_parameter_sensitivity.png}
      \vspace{0.1em}
      {\scriptsize 气泡大小 = TC 提升，颜色 = SAD 变化}
      
    \column{0.32\textwidth}
      \begin{block}{参数趋势}
        \begin{itemize}
          \setlength{\itemsep}{0.2em}
          \item 增大 sigma → TC 提升增加
          \item 增大 sigma → SAD 恶化加剧
          \item 增大窗口 → 效果饱和
        \end{itemize}
      \end{block}
      
      \vspace{0.3em}
      \begin{block}{推荐参数}
        \centering
        \metricchip{softgreen}{sigma=2.0}{window=5} \\
        \vspace{0.1em}
        {\scriptsize TC 提升: 4.0\%, SAD 增加: 2.4\%}
      \end{block}
  \end{columns}
  
  \vspace{0.1em}
  \notestrip{accentblue}{\scriptsize
  \textbf{建议：}追求平衡可选 Gaussian；追求精度无损可选 Bidirectional。
  }
\end{frame}

\begin{frame}[t]{综合性能雷达图}
  \vspace{-0.5cm}
  \scriptsize
  \linespread{0.98}\selectfont
  
  \begin{columns}[T,onlytextwidth]
    \column{0.60\textwidth}
      \centering
      \includegraphics[width=0.95\linewidth]{pic/figure5_radar_chart.png}
      
    \column{0.32\textwidth}
      \begin{block}{解读}
        \begin{itemize}
          \setlength{\itemsep}{0.2em}
          \item 雷达图中 SAD/MSE/Gradient 已反向归一化，\hilite{值越大 = 性能越好}。
          \item \textbf{Bidirectional} 最均衡
          \item Median 偏科严重
        \end{itemize}
      \end{block}
      
      \vspace{0.3em}
      \begin{block}{总结}
        \centering
        \textbf{Bidirectional (alpha=0.7)} \\
        \vspace{0.1em}
        → 雷达图面积最大 \\
        → 推荐为默认平滑策略
      \end{block}
  \end{columns}
  
  
\end{frame}

\begin{frame}[t]{视频换背景 Demo}
  \scriptsize
  \linespread{0.98}\selectfont
  
  \begin{columns}[T,onlytextwidth]
    \column{0.48\textwidth}
      \begin{block}{平滑后 alpha 对比}
        \centering
        \includegraphics[width=0.95\linewidth]{pic/gif_frames/frame_0001.png}
        \vspace{0.1em}
        {\scriptsize 左：原始 \quad | \quad 右：Bidirectional 平滑后}
        
        \vspace{0.2em}
        \href{run:pic/youtubematte_motion_0016_bidirectional_slow_hq.gif}{%
          \color{accentblue}{\underline{[Click] 播放动画效果 (GIF)}}%
        }
      \end{block}
      
    \column{0.48\textwidth}
      \begin{block}{换背景效果}
        \centering
        
\includegraphics[width=0.35\linewidth]{pic/background.jpg}
        \vspace{0.1em}
        \\
        {\scriptsize 平滑后 alpha → RGBA → 背景替换}
        
        \vspace{0.2em}
        \href{run:pic/youtubematte_motion_0016_bidirectional_composer.mp4}{%
          \color{accentblue}{\underline{[Click] 播放换背景动画 (mp4)}}%
        }
      \end{block}
  \end{columns}
  
  \vspace{0.2em}
  \begin{block}{接口说明}
    \scriptsize
    \textbf{输入：}图像主线首帧 alpha / mask
    \textbf{处理：}MatAnyone2 + Bidirectional 时序平滑 (alpha=0.7)
    \textbf{输出：}逐帧稳定 alpha + 换背景合成视频
  \end{block}
  
  \vspace{0.1em}
  \notestrip{softgreen}{\scriptsize
  \textbf{效果：}Bidirectional 平滑后，背景替换视频的\hilite{边界闪烁明显减少}，同时\hilite{发丝细节保持清晰}。
  }
\end{frame} 引入，
%    继承主文件导言区的全部宏 / 配色 / 字体，无需自带 \documentclass 或 \begin{document}。
%  · 直接在下面增删 \begin{frame}[t]{标题} ... \end{frame} 即可。
%  · 可复用样式：
%       \hilite{} \tcard{} \ccard{} \notestrip{} \metricchip{} \mpill[] \pill[]
%     可选 accent：nankai、accentblue、softgreen、softorange、softgray、deepred
% ============================================================================

\section{视频与应用}

\begin{frame}[t]{视频抠图与时序平滑}
  \scriptsize
  \linespread{0.98}\selectfont
  
  \begin{columns}[T,onlytextwidth]
    \column{0.48\textwidth}
      \begin{block}{问题：视频 alpha 抖动}
        \begin{itemize}
          \setlength{\itemsep}{0.2em}
          \item MatAnyone2~{\scriptsize[7]} 逐帧推理时，alpha 存在帧间\hilite{抖动与边界闪烁}；
          \item 直接替换背景会放大这些不稳定，导致视觉不自然；
          \item 需要一种\hilite{保持细节}的时序平滑策略。
        \end{itemize}
      \end{block}
      
      \vspace{0.2em}
      \begin{block}{解决方案：边界带时序平滑}
        \centering
        \resizebox{0.95\linewidth}{!}{%
        \begin{tikzpicture}[
          >={Latex[length=2.5mm,width=2mm]},
          node distance=3mm,
          every node/.style={font=\scriptsize},
          vcard/.style={rounded corners=2pt, line width=0.8pt, align=center, inner sep=3pt},
          vblue/.style={vcard, draw=accentblue!82!black, fill=accentblue!12},
          vpurple/.style={vcard, draw=nankai!88!black, fill=nankai!12},
          vgreen/.style={vcard, draw=softgreen!72!black, fill=softgreen!14},
          varrow/.style={-{Latex[length=2.5mm,width=2mm]}, line width=1.2pt, draw=nankai!80!black}
        ]
          \node[vblue] (input) {\textbf{输入}\\alpha video};
          \node[vpurple, right=of input] (detect) {\textbf{边界检测}\\$0.05<\alpha<0.95$};
          \node[vgreen, right=of detect] (smooth) {\textbf{时序平滑}\\EMA / Bidirectional};
          \node[vblue, right=of smooth] (output) {\textbf{输出}\\稳定 alpha};
          \draw[varrow] (input) -- (detect);
          \draw[varrow] (detect) -- (smooth);
          \draw[varrow] (smooth) -- (output);
        \end{tikzpicture}}
        
        \vspace{0.2em}
        \scriptsize
        仅在边界区域（$\alpha\in[0.05,0.95]$）施加平滑，\hilite{保护前景核区与背景}，避免过度模糊。
      \end{block}
      
    \column{0.48\textwidth}
      \begin{block}{7 种平滑策略对比}
        \centering
        \includegraphics[width=0.95\linewidth]{pic/figure6_performance_ranking.png}
        \vspace{0.1em}
        {\scriptsize TC 提升排名（左）与 SAD 变化排名（右）}
      \end{block}
  \end{columns}
\end{frame}

\begin{frame}[t]{7 种平滑策略性能评估}
  \scriptsize
  \linespread{0.50}\selectfont
  
  \begin{block}{实验设置}
    \scriptsize
    \begin{itemize}
      \setlength{\itemsep}{0.1em}
      \item \textbf{数据集：}YouTubeMatte (10 视频) + VideoMatt (10 视频)，共 20 个测试视频；
      \item \textbf{评估指标：}TC（时序一致性）、SAD（绝对误差和）、MSE、Gradient；
      \item \textbf{方法：}EMA、Gaussian、Median、Adaptive、Guided EMA、Bidirectional、Flow。
    \end{itemize}
  \end{block}
  
  \vspace{0.1em}
  \centering
  \resizebox{0.9\linewidth}{!}{%
  \begin{tabular}{lcccc}
    \toprule
    \textbf{方法} & \textbf{TC 提升 (\%)} & \textbf{SAD 变化 (\%)} & \textbf{MSE 变化 (\%)} & \textbf{细节保留} \\
    \midrule
    Median & 6.68 & +8.88 & +23.91 & Poor \\
    Gaussian & 5.50 & +5.40 & +11.42 & Fair \\
    Adaptive & 3.62 & +3.12 & +5.88 & Fair \\
    Guided EMA & 3.49 & +2.99 & +5.28 & Fair \\
    \textbf{Bidirectional} & \textbf{2.89} & \textbf{+0.44} & \textbf{-0.12} & \textbf{Excellent} \\
    EMA & 2.59 & +1.64 & +2.40 & Good \\
    Flow & 1.55 & +2.15 & +3.62 & Fair \\
    \bottomrule
  \end{tabular}}
  
  
  \vspace{0.1em}
  \notestrip{nankai}{\scriptsize
  \textbf{关键发现：}Bidirectional 是\hilite{唯一能同时改善 TC 和保持精度}的方法，细节保留最佳。
  }
\end{frame}

\begin{frame}[t]{方法对比散点图}
  \scriptsize
  \linespread{0.92}\selectfont
  
  \begin{columns}[T,onlytextwidth]
    \column{0.60\textwidth}
      \centering
      \includegraphics[width=\linewidth]{pic/figure1_tc_vs_sad_scatter.png}
      \vspace{0.1em}
      {\scriptsize  每个点代表一个视频的最佳结果}
      
    \column{0.4\textwidth}
      \begin{block}{解读}
        \begin{itemize}
          \setlength{\itemsep}{0.2em}
          \item 绿色区域：理想方向（TC↑，SAD↓）
          \item \textbf{Bidirectional}：最接近理想区域
          \item Median：TC 提升最大，但 SAD 增加显著
        \end{itemize}
      \end{block}
      
      \vspace{0.3em}
      \begin{block}{最佳参数}
        \scriptsize
        \begin{tabular}{ll}
          \toprule
          \textbf{方法} & \textbf{最佳参数} \\
          \midrule
          Median & window=7 \\
          Gaussian & sigma=2.0, window=7 \\
          \textbf{Bidirectional} & \textbf{alpha=0.7} \\
          \bottomrule
        \end{tabular}
      \end{block}
  \end{columns}
\end{frame}

\begin{frame}[t]{各方法性能分布箱线图}
  \scriptsize
  \linespread{0.98}\selectfont
  
  \begin{columns}[T,onlytextwidth]
    \column{0.7\textwidth}
      \centering
      \includegraphics[width=\linewidth]{pic/figure2_method_comparison_boxplot.png}
      \vspace{0.1em}
      {\scriptsize 左：TC 提升分布 \quad | \quad 右：SAD 变化分布}
      
    \column{0.32\textwidth}
      \begin{block}{箱线图解读}
        \begin{itemize}
          \setlength{\itemsep}{0.2em}
          \item 箱体 = 四分位距（50\% 数据）
          \item 横线 = 中位数
          \item 圆形 = 异常值
        \end{itemize}
      \end{block}
      
      \vspace{0.3em}
      \begin{block}{关键观察}
        \begin{itemize}
          \setlength{\itemsep}{0.2em}
          \item \textbf{Bidirectional}：SAD 分布最集中，异常值最少
          \item \textbf{Median}：TC 提升高但波动大
          \item \textbf{Gaussian}：TC 提升分布较稳定
        \end{itemize}
      \end{block}
  \end{columns}
  
  \vspace{0.1em}
  \notestrip{nankai}{\scriptsize
  \textbf{结论：}Bidirectional 的 SAD 变化分布最集中，说明其精度表现\hilite{最稳定}，不同视频间差异最小。
  }
\end{frame}

\begin{frame}[t]{Bidirectional 平滑效果对比}
  \scriptsize
  \linespread{0.98}\selectfont
  
  \begin{columns}[T,onlytextwidth]
    \column{0.55\textwidth}
      \centering
      \includegraphics[width=0.98\linewidth]{pic/gif_frames/frame_0001.png}
      \vspace{0.1em}
      {\tiny 左：原始 alpha \quad | \quad 右：Bidirectional 平滑后}
      
      \vspace{0.2em}
      \href{run:pic/youtubematte_motion_0016_bidirectional_comparison_slow_hq.gif}{%
        \color{accentblue}{\underline{[Click] 点击播放动画效果 (GIF)}}%
      }
      
      \vspace{0.3em}
      \begin{block}{视觉效果}
        \begin{itemize}
          \setlength{\itemsep}{0.2em}
          \item 发丝细节\hilite{无明显模糊}
          \item 边缘抖动\hilite{明显减弱}
          \item 背景替换后\hilite{闪烁减少}
        \end{itemize}
      \end{block}
      
    \column{0.42\textwidth}
      \begin{block}{性能指标}
        \scriptsize
        \begin{tabular}{lcc}
          \toprule
          \textbf{指标} & \textbf{数值} & \textbf{评价} \\
          \midrule
          TC 提升 & 2.89\% & 提升 \\
          SAD 变化 & +0.44\% & 几乎无损 \\
          MSE 变化 & -0.12\% & 改善 \\
          \bottomrule
        \end{tabular}
      \end{block}
      
      \vspace{0.3em}
      \begin{block}{优劣分析}
        \scriptsize
        \begin{itemize}
          \setlength{\itemsep}{0.1em}
          \item \textbf{优点：}精度几乎无损，发丝细节保留完整
          \item \textbf{缺点：}TC 提升幅度相对较小
        \end{itemize}
      \end{block}
  \end{columns}
  
  \vspace{0.2em}
  \notestrip{softgreen}{\scriptsize
  \textbf{结论：}Bidirectional 在\hilite{保持视觉质量的同时}实现时序一致性提升，是精度无损场景的最佳选择。
  }
\end{frame}

\begin{frame}[t]{Median 平滑效果对比}
  \scriptsize
  \linespread{0.98}\selectfont
  
  \begin{columns}[T,onlytextwidth]
    \column{0.55\textwidth}
      \centering
      \includegraphics[width=0.98\linewidth]{pic/median_frames/median_frame_0001.png}
      \vspace{0.1em}
      {\tiny 左：原始 alpha \quad | \quad 右：Median 平滑后}
      
      \vspace{0.2em}
      \href{run:pic/youtubematte_motion_0016_comparison_slow_hq.gif}{%
        \color{accentblue}{\underline{[Click] 点击播放动画效果 (GIF)}}%
      }
      
      \vspace{0.3em}
      \begin{block}{视觉效果}
        \begin{itemize}
          \setlength{\itemsep}{0.2em}
          \item 发丝细节\hilite{明显模糊}
          \item 边缘抖动\hilite{减少但细节丢失}
          \item 背景替换后\hilite{毛边明显}
        \end{itemize}
      \end{block}
      
    \column{0.42\textwidth}
      \begin{block}{性能指标}
        \scriptsize
        \begin{tabular}{lcc}
          \toprule
          \textbf{指标} & \textbf{数值} & \textbf{评价} \\
          \midrule
          TC 提升 & 6.68\% & 提升最大 \\
          SAD 变化 & +8.88\% & 恶化严重 \\
          MSE 变化 & +23.91\% & 恶化严重 \\
          \bottomrule
        \end{tabular}
      \end{block}
      
      \vspace{0.3em}
      \begin{block}{优劣分析}
        \scriptsize
        \begin{itemize}
          \setlength{\itemsep}{0.1em}
          \item \textbf{优点：}TC 提升最大（6.68\%），时序稳定性最好
          \item \textbf{缺点：}发丝模糊严重，视觉损失不可接受
        \end{itemize}
      \end{block}
  \end{columns}
  
  \vspace{0.2em}
  \notestrip{softorange}{\scriptsize
  \textbf{结论：}Median 数值最优但\hilite{视觉损失严重}，不适合需要精细边缘的实际应用。
  }
\end{frame}


\begin{frame}[t]{Gaussian 平滑效果对比}
  \scriptsize
  \linespread{0.98}\selectfont
  
  \begin{columns}[T,onlytextwidth]
    \column{0.55\textwidth}
      \centering
      
\includegraphics[width=0.98\linewidth]{pic/gaussian_frames/gaussian_frame_0001.png}
      \vspace{0.1em}
      {\tiny 左：原始 alpha \quad | \quad 右：Gaussian 平滑后 (sigma=2.0)}
      
      \vspace{0.2em}
      \href{run:pic/youtubematte_motion_0016_gaussian_comparison_slow_hq.gif}{%
        \color{accentblue}{\underline{[Click] 点击播放动画效果 (GIF)}}%
      }
      
      \vspace{0.3em}
      \begin{block}{视觉效果}
        \begin{itemize}
          \setlength{\itemsep}{0.2em}
          \item 发丝细节\hilite{轻微模糊}
          \item 边缘抖动\hilite{明显减少}
          \item 整体\hilite{平衡性较好}
        \end{itemize}
      \end{block}
      
    \column{0.42\textwidth}
      \begin{block}{性能指标}
        \scriptsize
        \begin{tabular}{lcc}
          \toprule
          \textbf{指标} & \textbf{数值} & \textbf{评价} \\
          \midrule
          TC 提升 & 4.0\% & 提升显著 \\
          SAD 变化 & +2.4\% & 轻微恶化 \\
          MSE 变化 & +3.0\% & 轻微恶化 \\
          \bottomrule
        \end{tabular}
      \end{block}
      
      \vspace{0.3em}
      \begin{block}{优劣分析}
        \scriptsize
        \begin{itemize}
          \setlength{\itemsep}{0.1em}
          \item \textbf{优点：}TC 提升显著（4.0\%），是 TC 和精度的最佳平衡点
          \item \textbf{缺点：}发丝有轻微模糊，精度略有下降
        \end{itemize}
      \end{block}
  \end{columns}
  
  \vspace{0.2em}
  \notestrip{accentblue}{\scriptsize
  \textbf{建议：}追求平衡可选 Gaussian (sigma=2.0, window=5)；追求精度无损可选 Bidirectional (alpha=0.7)。
  }
\end{frame}


\begin{frame}[t]{两个数据集上的性能对比}
  \scriptsize
  \linespread{0.98}\selectfont
  
  \begin{columns}[T,onlytextwidth]
    \column{0.7\textwidth}
      \centering
      \includegraphics[width=0.98\linewidth]{pic/figure3_dataset_comparison.png}
      \vspace{0.1em}
      {\scriptsize 左：TC 提升对比 \quad | \quad 右：SAD 变化对比}
      
    \column{0.32\textwidth}
      \begin{block}{关键观察}
        \begin{itemize}
          \setlength{\itemsep}{0.2em}
          \item YouTubeMatte 上效果\hilite{普遍更好}
          \item VideoMatt 更具挑战性
          \item \textbf{Bidirectional} 表现\hilite{最稳定}
        \end{itemize}
      \end{block}
      
      \vspace{0.5em}
      \begin{block}{差异分析}
        \scriptsize
        VideoMatt 包含更多\hilite{复杂运动}和\hilite{精细结构}，平滑难度更高。
      \end{block}
    \end{columns}
   \vspace{0.1em}
  \notestrip{accentblue}{\scriptsize
  \textbf{结论：}Bidirectional 在两个数据集上表现最稳定，适合不同场景的视频抠图任务。
  }
\end{frame}

\begin{frame}[t]{Gaussian 参数敏感性分析}
  \vspace{-0.5cm}
  \scriptsize
  \linespread{0.98}\selectfont
  
  \begin{columns}[T,onlytextwidth]
    \column{0.65\textwidth}
      \centering
      \includegraphics[width=0.98\linewidth]{pic/figure4_gaussian_parameter_sensitivity.png}
      \vspace{0.1em}
      {\scriptsize 气泡大小 = TC 提升，颜色 = SAD 变化}
      
    \column{0.32\textwidth}
      \begin{block}{参数趋势}
        \begin{itemize}
          \setlength{\itemsep}{0.2em}
          \item 增大 sigma → TC 提升增加
          \item 增大 sigma → SAD 恶化加剧
          \item 增大窗口 → 效果饱和
        \end{itemize}
      \end{block}
      
      \vspace{0.3em}
      \begin{block}{推荐参数}
        \centering
        \metricchip{softgreen}{sigma=2.0}{window=5} \\
        \vspace{0.1em}
        {\scriptsize TC 提升: 4.0\%, SAD 增加: 2.4\%}
      \end{block}
  \end{columns}
  
  \vspace{0.1em}
  \notestrip{accentblue}{\scriptsize
  \textbf{建议：}追求平衡可选 Gaussian；追求精度无损可选 Bidirectional。
  }
\end{frame}

\begin{frame}[t]{综合性能雷达图}
  \vspace{-0.5cm}
  \scriptsize
  \linespread{0.98}\selectfont
  
  \begin{columns}[T,onlytextwidth]
    \column{0.60\textwidth}
      \centering
      \includegraphics[width=0.95\linewidth]{pic/figure5_radar_chart.png}
      
    \column{0.32\textwidth}
      \begin{block}{解读}
        \begin{itemize}
          \setlength{\itemsep}{0.2em}
          \item 雷达图中 SAD/MSE/Gradient 已反向归一化，\hilite{值越大 = 性能越好}。
          \item \textbf{Bidirectional} 最均衡
          \item Median 偏科严重
        \end{itemize}
      \end{block}
      
      \vspace{0.3em}
      \begin{block}{总结}
        \centering
        \textbf{Bidirectional (alpha=0.7)} \\
        \vspace{0.1em}
        → 雷达图面积最大 \\
        → 推荐为默认平滑策略
      \end{block}
  \end{columns}
  
  
\end{frame}

\begin{frame}[t]{视频换背景 Demo}
  \scriptsize
  \linespread{0.98}\selectfont
  
  \begin{columns}[T,onlytextwidth]
    \column{0.48\textwidth}
      \begin{block}{平滑后 alpha 对比}
        \centering
        \includegraphics[width=0.95\linewidth]{pic/gif_frames/frame_0001.png}
        \vspace{0.1em}
        {\scriptsize 左：原始 \quad | \quad 右：Bidirectional 平滑后}
        
        \vspace{0.2em}
        \href{run:pic/youtubematte_motion_0016_bidirectional_slow_hq.gif}{%
          \color{accentblue}{\underline{[Click] 播放动画效果 (GIF)}}%
        }
      \end{block}
      
    \column{0.48\textwidth}
      \begin{block}{换背景效果}
        \centering
        
\includegraphics[width=0.35\linewidth]{pic/background.jpg}
        \vspace{0.1em}
        \\
        {\scriptsize 平滑后 alpha → RGBA → 背景替换}
        
        \vspace{0.2em}
        \href{run:pic/youtubematte_motion_0016_bidirectional_composer.mp4}{%
          \color{accentblue}{\underline{[Click] 播放换背景动画 (mp4)}}%
        }
      \end{block}
  \end{columns}
  
  \vspace{0.2em}
  \begin{block}{接口说明}
    \scriptsize
    \textbf{输入：}图像主线首帧 alpha / mask
    \textbf{处理：}MatAnyone2 + Bidirectional 时序平滑 (alpha=0.7)
    \textbf{输出：}逐帧稳定 alpha + 换背景合成视频
  \end{block}
  
  \vspace{0.1em}
  \notestrip{softgreen}{\scriptsize
  \textbf{效果：}Bidirectional 平滑后，背景替换视频的\hilite{边界闪烁明显减少}，同时\hilite{发丝细节保持清晰}。
  }
\end{frame} 引入，
%    继承主文件导言区的全部宏 / 配色 / 字体，无需自带 \documentclass 或 \begin{document}。
%  · 直接在下面增删 \begin{frame}[t]{标题} ... \end{frame} 即可。
%  · 可复用样式：
%       \hilite{} \tcard{} \ccard{} \notestrip{} \metricchip{} \mpill[] \pill[]
%     可选 accent：nankai、accentblue、softgreen、softorange、softgray、deepred
% ============================================================================

\section{视频与应用}

\begin{frame}[t]{视频抠图与时序平滑}
  \scriptsize
  \linespread{0.98}\selectfont
  
  \begin{columns}[T,onlytextwidth]
    \column{0.48\textwidth}
      \begin{block}{问题：视频 alpha 抖动}
        \begin{itemize}
          \setlength{\itemsep}{0.2em}
          \item MatAnyone2~{\scriptsize[7]} 逐帧推理时，alpha 存在帧间\hilite{抖动与边界闪烁}；
          \item 直接替换背景会放大这些不稳定，导致视觉不自然；
          \item 需要一种\hilite{保持细节}的时序平滑策略。
        \end{itemize}
      \end{block}
      
      \vspace{0.2em}
      \begin{block}{解决方案：边界带时序平滑}
        \centering
        \resizebox{0.95\linewidth}{!}{%
        \begin{tikzpicture}[
          >={Latex[length=2.5mm,width=2mm]},
          node distance=3mm,
          every node/.style={font=\scriptsize},
          vcard/.style={rounded corners=2pt, line width=0.8pt, align=center, inner sep=3pt},
          vblue/.style={vcard, draw=accentblue!82!black, fill=accentblue!12},
          vpurple/.style={vcard, draw=nankai!88!black, fill=nankai!12},
          vgreen/.style={vcard, draw=softgreen!72!black, fill=softgreen!14},
          varrow/.style={-{Latex[length=2.5mm,width=2mm]}, line width=1.2pt, draw=nankai!80!black}
        ]
          \node[vblue] (input) {\textbf{输入}\\alpha video};
          \node[vpurple, right=of input] (detect) {\textbf{边界检测}\\$0.05<\alpha<0.95$};
          \node[vgreen, right=of detect] (smooth) {\textbf{时序平滑}\\EMA / Bidirectional};
          \node[vblue, right=of smooth] (output) {\textbf{输出}\\稳定 alpha};
          \draw[varrow] (input) -- (detect);
          \draw[varrow] (detect) -- (smooth);
          \draw[varrow] (smooth) -- (output);
        \end{tikzpicture}}
        
        \vspace{0.2em}
        \scriptsize
        仅在边界区域（$\alpha\in[0.05,0.95]$）施加平滑，\hilite{保护前景核区与背景}，避免过度模糊。
      \end{block}
      
    \column{0.48\textwidth}
      \begin{block}{7 种平滑策略对比}
        \centering
        \includegraphics[width=0.95\linewidth]{pic/figure6_performance_ranking.png}
        \vspace{0.1em}
        {\scriptsize TC 提升排名（左）与 SAD 变化排名（右）}
      \end{block}
  \end{columns}
\end{frame}

\begin{frame}[t]{7 种平滑策略性能评估}
  \scriptsize
  \linespread{0.50}\selectfont
  
  \begin{block}{实验设置}
    \scriptsize
    \begin{itemize}
      \setlength{\itemsep}{0.1em}
      \item \textbf{数据集：}YouTubeMatte (10 视频) + VideoMatt (10 视频)，共 20 个测试视频；
      \item \textbf{评估指标：}TC（时序一致性）、SAD（绝对误差和）、MSE、Gradient；
      \item \textbf{方法：}EMA、Gaussian、Median、Adaptive、Guided EMA、Bidirectional、Flow。
    \end{itemize}
  \end{block}
  
  \vspace{0.1em}
  \centering
  \resizebox{0.9\linewidth}{!}{%
  \begin{tabular}{lcccc}
    \toprule
    \textbf{方法} & \textbf{TC 提升 (\%)} & \textbf{SAD 变化 (\%)} & \textbf{MSE 变化 (\%)} & \textbf{细节保留} \\
    \midrule
    Median & 6.68 & +8.88 & +23.91 & Poor \\
    Gaussian & 5.50 & +5.40 & +11.42 & Fair \\
    Adaptive & 3.62 & +3.12 & +5.88 & Fair \\
    Guided EMA & 3.49 & +2.99 & +5.28 & Fair \\
    \textbf{Bidirectional} & \textbf{2.89} & \textbf{+0.44} & \textbf{-0.12} & \textbf{Excellent} \\
    EMA & 2.59 & +1.64 & +2.40 & Good \\
    Flow & 1.55 & +2.15 & +3.62 & Fair \\
    \bottomrule
  \end{tabular}}
  
  
  \vspace{0.1em}
  \notestrip{nankai}{\scriptsize
  \textbf{关键发现：}Bidirectional 是\hilite{唯一能同时改善 TC 和保持精度}的方法，细节保留最佳。
  }
\end{frame}

\begin{frame}[t]{方法对比散点图}
  \scriptsize
  \linespread{0.92}\selectfont
  
  \begin{columns}[T,onlytextwidth]
    \column{0.60\textwidth}
      \centering
      \includegraphics[width=\linewidth]{pic/figure1_tc_vs_sad_scatter.png}
      \vspace{0.1em}
      {\scriptsize  每个点代表一个视频的最佳结果}
      
    \column{0.4\textwidth}
      \begin{block}{解读}
        \begin{itemize}
          \setlength{\itemsep}{0.2em}
          \item 绿色区域：理想方向（TC↑，SAD↓）
          \item \textbf{Bidirectional}：最接近理想区域
          \item Median：TC 提升最大，但 SAD 增加显著
        \end{itemize}
      \end{block}
      
      \vspace{0.3em}
      \begin{block}{最佳参数}
        \scriptsize
        \begin{tabular}{ll}
          \toprule
          \textbf{方法} & \textbf{最佳参数} \\
          \midrule
          Median & window=7 \\
          Gaussian & sigma=2.0, window=7 \\
          \textbf{Bidirectional} & \textbf{alpha=0.7} \\
          \bottomrule
        \end{tabular}
      \end{block}
  \end{columns}
\end{frame}

\begin{frame}[t]{各方法性能分布箱线图}
  \scriptsize
  \linespread{0.98}\selectfont
  
  \begin{columns}[T,onlytextwidth]
    \column{0.7\textwidth}
      \centering
      \includegraphics[width=\linewidth]{pic/figure2_method_comparison_boxplot.png}
      \vspace{0.1em}
      {\scriptsize 左：TC 提升分布 \quad | \quad 右：SAD 变化分布}
      
    \column{0.32\textwidth}
      \begin{block}{箱线图解读}
        \begin{itemize}
          \setlength{\itemsep}{0.2em}
          \item 箱体 = 四分位距（50\% 数据）
          \item 横线 = 中位数
          \item 圆形 = 异常值
        \end{itemize}
      \end{block}
      
      \vspace{0.3em}
      \begin{block}{关键观察}
        \begin{itemize}
          \setlength{\itemsep}{0.2em}
          \item \textbf{Bidirectional}：SAD 分布最集中，异常值最少
          \item \textbf{Median}：TC 提升高但波动大
          \item \textbf{Gaussian}：TC 提升分布较稳定
        \end{itemize}
      \end{block}
  \end{columns}
  
  \vspace{0.1em}
  \notestrip{nankai}{\scriptsize
  \textbf{结论：}Bidirectional 的 SAD 变化分布最集中，说明其精度表现\hilite{最稳定}，不同视频间差异最小。
  }
\end{frame}

\begin{frame}[t]{Bidirectional 平滑效果对比}
  \scriptsize
  \linespread{0.98}\selectfont
  
  \begin{columns}[T,onlytextwidth]
    \column{0.55\textwidth}
      \centering
      \includegraphics[width=0.98\linewidth]{pic/gif_frames/frame_0001.png}
      \vspace{0.1em}
      {\tiny 左：原始 alpha \quad | \quad 右：Bidirectional 平滑后}
      
      \vspace{0.2em}
      \href{run:pic/youtubematte_motion_0016_bidirectional_comparison_slow_hq.gif}{%
        \color{accentblue}{\underline{[Click] 点击播放动画效果 (GIF)}}%
      }
      
      \vspace{0.3em}
      \begin{block}{视觉效果}
        \begin{itemize}
          \setlength{\itemsep}{0.2em}
          \item 发丝细节\hilite{无明显模糊}
          \item 边缘抖动\hilite{明显减弱}
          \item 背景替换后\hilite{闪烁减少}
        \end{itemize}
      \end{block}
      
    \column{0.42\textwidth}
      \begin{block}{性能指标}
        \scriptsize
        \begin{tabular}{lcc}
          \toprule
          \textbf{指标} & \textbf{数值} & \textbf{评价} \\
          \midrule
          TC 提升 & 2.89\% & 提升 \\
          SAD 变化 & +0.44\% & 几乎无损 \\
          MSE 变化 & -0.12\% & 改善 \\
          \bottomrule
        \end{tabular}
      \end{block}
      
      \vspace{0.3em}
      \begin{block}{优劣分析}
        \scriptsize
        \begin{itemize}
          \setlength{\itemsep}{0.1em}
          \item \textbf{优点：}精度几乎无损，发丝细节保留完整
          \item \textbf{缺点：}TC 提升幅度相对较小
        \end{itemize}
      \end{block}
  \end{columns}
  
  \vspace{0.2em}
  \notestrip{softgreen}{\scriptsize
  \textbf{结论：}Bidirectional 在\hilite{保持视觉质量的同时}实现时序一致性提升，是精度无损场景的最佳选择。
  }
\end{frame}

\begin{frame}[t]{Median 平滑效果对比}
  \scriptsize
  \linespread{0.98}\selectfont
  
  \begin{columns}[T,onlytextwidth]
    \column{0.55\textwidth}
      \centering
      \includegraphics[width=0.98\linewidth]{pic/median_frames/median_frame_0001.png}
      \vspace{0.1em}
      {\tiny 左：原始 alpha \quad | \quad 右：Median 平滑后}
      
      \vspace{0.2em}
      \href{run:pic/youtubematte_motion_0016_comparison_slow_hq.gif}{%
        \color{accentblue}{\underline{[Click] 点击播放动画效果 (GIF)}}%
      }
      
      \vspace{0.3em}
      \begin{block}{视觉效果}
        \begin{itemize}
          \setlength{\itemsep}{0.2em}
          \item 发丝细节\hilite{明显模糊}
          \item 边缘抖动\hilite{减少但细节丢失}
          \item 背景替换后\hilite{毛边明显}
        \end{itemize}
      \end{block}
      
    \column{0.42\textwidth}
      \begin{block}{性能指标}
        \scriptsize
        \begin{tabular}{lcc}
          \toprule
          \textbf{指标} & \textbf{数值} & \textbf{评价} \\
          \midrule
          TC 提升 & 6.68\% & 提升最大 \\
          SAD 变化 & +8.88\% & 恶化严重 \\
          MSE 变化 & +23.91\% & 恶化严重 \\
          \bottomrule
        \end{tabular}
      \end{block}
      
      \vspace{0.3em}
      \begin{block}{优劣分析}
        \scriptsize
        \begin{itemize}
          \setlength{\itemsep}{0.1em}
          \item \textbf{优点：}TC 提升最大（6.68\%），时序稳定性最好
          \item \textbf{缺点：}发丝模糊严重，视觉损失不可接受
        \end{itemize}
      \end{block}
  \end{columns}
  
  \vspace{0.2em}
  \notestrip{softorange}{\scriptsize
  \textbf{结论：}Median 数值最优但\hilite{视觉损失严重}，不适合需要精细边缘的实际应用。
  }
\end{frame}


\begin{frame}[t]{Gaussian 平滑效果对比}
  \scriptsize
  \linespread{0.98}\selectfont
  
  \begin{columns}[T,onlytextwidth]
    \column{0.55\textwidth}
      \centering
      
\includegraphics[width=0.98\linewidth]{pic/gaussian_frames/gaussian_frame_0001.png}
      \vspace{0.1em}
      {\tiny 左：原始 alpha \quad | \quad 右：Gaussian 平滑后 (sigma=2.0)}
      
      \vspace{0.2em}
      \href{run:pic/youtubematte_motion_0016_gaussian_comparison_slow_hq.gif}{%
        \color{accentblue}{\underline{[Click] 点击播放动画效果 (GIF)}}%
      }
      
      \vspace{0.3em}
      \begin{block}{视觉效果}
        \begin{itemize}
          \setlength{\itemsep}{0.2em}
          \item 发丝细节\hilite{轻微模糊}
          \item 边缘抖动\hilite{明显减少}
          \item 整体\hilite{平衡性较好}
        \end{itemize}
      \end{block}
      
    \column{0.42\textwidth}
      \begin{block}{性能指标}
        \scriptsize
        \begin{tabular}{lcc}
          \toprule
          \textbf{指标} & \textbf{数值} & \textbf{评价} \\
          \midrule
          TC 提升 & 4.0\% & 提升显著 \\
          SAD 变化 & +2.4\% & 轻微恶化 \\
          MSE 变化 & +3.0\% & 轻微恶化 \\
          \bottomrule
        \end{tabular}
      \end{block}
      
      \vspace{0.3em}
      \begin{block}{优劣分析}
        \scriptsize
        \begin{itemize}
          \setlength{\itemsep}{0.1em}
          \item \textbf{优点：}TC 提升显著（4.0\%），是 TC 和精度的最佳平衡点
          \item \textbf{缺点：}发丝有轻微模糊，精度略有下降
        \end{itemize}
      \end{block}
  \end{columns}
  
  \vspace{0.2em}
  \notestrip{accentblue}{\scriptsize
  \textbf{建议：}追求平衡可选 Gaussian (sigma=2.0, window=5)；追求精度无损可选 Bidirectional (alpha=0.7)。
  }
\end{frame}


\begin{frame}[t]{两个数据集上的性能对比}
  \scriptsize
  \linespread{0.98}\selectfont
  
  \begin{columns}[T,onlytextwidth]
    \column{0.7\textwidth}
      \centering
      \includegraphics[width=0.98\linewidth]{pic/figure3_dataset_comparison.png}
      \vspace{0.1em}
      {\scriptsize 左：TC 提升对比 \quad | \quad 右：SAD 变化对比}
      
    \column{0.32\textwidth}
      \begin{block}{关键观察}
        \begin{itemize}
          \setlength{\itemsep}{0.2em}
          \item YouTubeMatte 上效果\hilite{普遍更好}
          \item VideoMatt 更具挑战性
          \item \textbf{Bidirectional} 表现\hilite{最稳定}
        \end{itemize}
      \end{block}
      
      \vspace{0.5em}
      \begin{block}{差异分析}
        \scriptsize
        VideoMatt 包含更多\hilite{复杂运动}和\hilite{精细结构}，平滑难度更高。
      \end{block}
    \end{columns}
   \vspace{0.1em}
  \notestrip{accentblue}{\scriptsize
  \textbf{结论：}Bidirectional 在两个数据集上表现最稳定，适合不同场景的视频抠图任务。
  }
\end{frame}

\begin{frame}[t]{Gaussian 参数敏感性分析}
  \vspace{-0.5cm}
  \scriptsize
  \linespread{0.98}\selectfont
  
  \begin{columns}[T,onlytextwidth]
    \column{0.65\textwidth}
      \centering
      \includegraphics[width=0.98\linewidth]{pic/figure4_gaussian_parameter_sensitivity.png}
      \vspace{0.1em}
      {\scriptsize 气泡大小 = TC 提升，颜色 = SAD 变化}
      
    \column{0.32\textwidth}
      \begin{block}{参数趋势}
        \begin{itemize}
          \setlength{\itemsep}{0.2em}
          \item 增大 sigma → TC 提升增加
          \item 增大 sigma → SAD 恶化加剧
          \item 增大窗口 → 效果饱和
        \end{itemize}
      \end{block}
      
      \vspace{0.3em}
      \begin{block}{推荐参数}
        \centering
        \metricchip{softgreen}{sigma=2.0}{window=5} \\
        \vspace{0.1em}
        {\scriptsize TC 提升: 4.0\%, SAD 增加: 2.4\%}
      \end{block}
  \end{columns}
  
  \vspace{0.1em}
  \notestrip{accentblue}{\scriptsize
  \textbf{建议：}追求平衡可选 Gaussian；追求精度无损可选 Bidirectional。
  }
\end{frame}

\begin{frame}[t]{综合性能雷达图}
  \vspace{-0.5cm}
  \scriptsize
  \linespread{0.98}\selectfont
  
  \begin{columns}[T,onlytextwidth]
    \column{0.60\textwidth}
      \centering
      \includegraphics[width=0.95\linewidth]{pic/figure5_radar_chart.png}
      
    \column{0.32\textwidth}
      \begin{block}{解读}
        \begin{itemize}
          \setlength{\itemsep}{0.2em}
          \item 雷达图中 SAD/MSE/Gradient 已反向归一化，\hilite{值越大 = 性能越好}。
          \item \textbf{Bidirectional} 最均衡
          \item Median 偏科严重
        \end{itemize}
      \end{block}
      
      \vspace{0.3em}
      \begin{block}{总结}
        \centering
        \textbf{Bidirectional (alpha=0.7)} \\
        \vspace{0.1em}
        → 雷达图面积最大 \\
        → 推荐为默认平滑策略
      \end{block}
  \end{columns}
  
  
\end{frame}

\begin{frame}[t]{视频换背景 Demo}
  \scriptsize
  \linespread{0.98}\selectfont
  
  \begin{columns}[T,onlytextwidth]
    \column{0.48\textwidth}
      \begin{block}{平滑后 alpha 对比}
        \centering
        \includegraphics[width=0.95\linewidth]{pic/gif_frames/frame_0001.png}
        \vspace{0.1em}
        {\scriptsize 左：原始 \quad | \quad 右：Bidirectional 平滑后}
        
        \vspace{0.2em}
        \href{run:pic/youtubematte_motion_0016_bidirectional_slow_hq.gif}{%
          \color{accentblue}{\underline{[Click] 播放动画效果 (GIF)}}%
        }
      \end{block}
      
    \column{0.48\textwidth}
      \begin{block}{换背景效果}
        \centering
        
\includegraphics[width=0.35\linewidth]{pic/background.jpg}
        \vspace{0.1em}
        \\
        {\scriptsize 平滑后 alpha → RGBA → 背景替换}
        
        \vspace{0.2em}
        \href{run:pic/youtubematte_motion_0016_bidirectional_composer.mp4}{%
          \color{accentblue}{\underline{[Click] 播放换背景动画 (mp4)}}%
        }
      \end{block}
  \end{columns}
  
  \vspace{0.2em}
  \begin{block}{接口说明}
    \scriptsize
    \textbf{输入：}图像主线首帧 alpha / mask
    \textbf{处理：}MatAnyone2 + Bidirectional 时序平滑 (alpha=0.7)
    \textbf{输出：}逐帧稳定 alpha + 换背景合成视频
  \end{block}
  
  \vspace{0.1em}
  \notestrip{softgreen}{\scriptsize
  \textbf{效果：}Bidirectional 平滑后，背景替换视频的\hilite{边界闪烁明显减少}，同时\hilite{发丝细节保持清晰}。
  }
\end{frame}